\section{Diskussion}
\label{sec:Diskussion}

Die relative Messabweichung wurde nach der Formel 
\begin{equation}
    \Delta x =   \frac{ \left| x_{exp} - x_{theo} \right|   }{x_{theo}}  
    \label{eq:abweichung}
\end{equation}
bestimmt.

\subsection{Bestimmung des B-Feldes}
Die Bestimmung des B-Feldes hat gut funktioniert, hierbei sind keine offensichtlichen Fehler aufgefallen.
\subsection{Bestimmung der effektiven Masse}
Um einen Vergleich zu ziehen, ob die Messung gelungen ist wird der Wert der effektiven Masse mit dem Theoriewert von GaAs verglichen. Hierzu folgt der Theoriewert von GaAs \cite{GaAs} mit 
\begin{equation*}
    m^\ast = 0.067 \, m_e.
\end{equation*}
Daraus ergibt sich nach \eqref{eq:abweichung} eine relative Messabweichung von 
\begin{align*}
    m_{\text{eff,1}} &= 31.34\,\% \\
    m_{\text{eff,2}} &= 41.79\, \%
\end{align*}
Der Fehler kann hierbei zustande kommen, wegen ungenauem Ablesen der Messinstrumente oder es kann der Polarisationfilter beispielsweise nicht perfekt eingestellt sein für die einzelnen Messungen. Ein weiterer Grund für den Fehler kann die Wärme des Magneten gewesen sein, welcher im Verlauf des Versuchs immer heißer geworden ist, wodurch die Leitfähigkeit des Halbleiters erhöht wird. 
Aufgrund von beispielsweise veralteter und fehlerbehafteter Messgeräte kann es des weiteren zu Fehlern kommen, so musste beispielsweise ein Messwert in \autoref{fig:masse2} in der linearen Regression ignoriert werden.
Was durch die Messung aber bestätigt werden konnte ist, dass die dichter dotierte Probe eine höhere effektive Masse hat als die weniger dicht dotierte Probe, was dadurch erklärt werden kann, dass bei GaAs die Bandstruktur einen großen Anharmonischen Anteil hat, sich die effektive Masse also stärker mit der Energie der Elektronen ändert und somit bei einer erhöhten Dichte von Leitungselektronen auch stärker von der weniger stark dotierten Probe abweicht.
\clearpage